%%%%%%%%%%%%%%%%%%%%%%%%%%%%%%%%%%%%%%%%%%%%%%%%%%%%%%%%%%%%%%%%%%%%%%
%% Springer Nature LaTeX template (sn-jnl) — npj Digital Medicine submission
%% Converted from the Elsevier (elsarticle) source paper_elsevier.tex
%% Submit as ONE .tex document; figures attached separately (do not \input).
%%%%%%%%%%%%%%%%%%%%%%%%%%%%%%%%%%%%%%%%%%%%%%%%%%%%%%%%%%%%%%%%%%%%%%
\documentclass[pdflatex,sn-nature]{sn-jnl}% Nature Portfolio Reference Style

%%%% Standard packages
\usepackage{graphicx}%
\usepackage{multirow}%
\usepackage{amsmath,amssymb,amsfonts}%
\usepackage{amsthm}%
\usepackage{mathrsfs}%
\usepackage[title]{appendix}%
\usepackage{xcolor}%
\usepackage{textcomp}%
\usepackage{manyfoot}%
\usepackage{booktabs}%
\usepackage{tabularx}%
\usepackage{algorithm}%
\usepackage{algorithmic}%
\usepackage{url}%

%% Math operators and macros carried over from the source manuscript
\DeclareMathOperator*{\argmin}{arg\,min}
\DeclareMathOperator{\Var}{Var}
\DeclareMathOperator{\Tr}{Tr}
\newcommand{\vect}[1]{\mathbf{#1}}
\newcommand{\matr}[1]{\mathbf{#1}}
\newcommand{\set}[1]{\mathcal{#1}}

%% Theorem-like environments (template styles)
\theoremstyle{thmstyleone}%
\newtheorem{theorem}{Theorem}%
\newtheorem{proposition}[theorem]{Proposition}%

\theoremstyle{thmstyletwo}%
\newtheorem{example}{Example}%
\newtheorem{remark}{Remark}%

\theoremstyle{thmstylethree}%
\newtheorem{definition}{Definition}%

\raggedbottom

\begin{document}

\title[Deep Conformal Alignment for Histopathology]{Deep Conformal Alignment: Learning Canonical Representations for Robust Histopathology Image Analysis}

%%=============================================================%%
%% Replace the placeholder author/affiliation block below with
%% the real author list before submission.
%%=============================================================%%
\author*[1]{\fnm{First} \sur{Author}}\email{author1@example.com}

\author[2]{\fnm{Second} \sur{Author}}\email{author2@example.com}

\affil*[1]{\orgdiv{Department}, \orgname{Organization 1}, \orgaddress{\city{City}, \country{Country}}}

\affil[2]{\orgdiv{Department}, \orgname{Organization 2}, \orgaddress{\city{City}, \country{Country}}}

%%==================================%%
%% Unstructured abstract            %%
%%==================================%%
\abstract{Histopathology image analysis faces critical challenges from morphological variability due to tissue preparation artifacts, staining inconsistencies, and biological heterogeneity, requiring large annotated datasets that are expensive to obtain. While recent geometric normalization methods using fixed conformal mappings improve data efficiency, they require per-image optimization and cannot adapt to tissue semantics. We propose Deep Conformal Alignment (DCA), a principled framework integrating conformal geometry into end-to-end neural network learning. DCA employs a diffeomorphic spatial transformer parameterized by velocity fields with scaling-and-squaring integration, ensuring topology preservation. Our key contribution is a differentiable Conformal Energy Loss derived from the Cauchy-Riemann equations, connected to quasiconformal mapping theory through the Beltrami coefficient, which guides networks toward angle-preserving transformations that maintain local tissue structure. Comprehensive evaluation on colorectal cancer histopathology with 100,000 training images demonstrates that DCA achieves 93.2\% accuracy across 9 tissue types, outperforming diverse baselines including spatial transformers, fixed geometric normalization, and self-supervised learning. DCA is markedly label-efficient: trained on 10\% of labeled data (66.4\% accuracy) it matches supervised baselines trained on several times more annotations (reported peak label-efficiency up to 4.2$\times$). While fine-tuned foundation models achieve strong performance (94.3\%), DCA complements rather than competes with them: combining DCA's geometric normalization with foundation model representations achieves 95.1\% accuracy, demonstrating synergistic benefits. Learned transformations are semantically adaptive: tumor tissue undergoes aggressive normalization while structured tissues receive subtle adjustments, with deformation magnitude correlating with tissue morphological variability (r=0.94). Our work establishes principled integration of differential geometry into deep learning, offering a promising direction for data-efficient medical image analysis that remains valuable in the foundation model era.}

\keywords{Histopathology, Conformal geometry, Deep learning, Image classification, Data efficiency}

\maketitle
